
(G) \uline{Il
existe un} $Y$-\uline{schéma lisse} $V$, \uline{factorisant} $f$, \uline{tel que pour toute immersion fermée} : $Y\to V'$ \uline{on puisse trouver un} $V'$-\uline{schéma lisse} $V''$, \uline{avec} $V\simeq V''\times_{V'}Y$ (\uline{cette condition est en particulier remplie si} $f$ \uline{est quasi-projectif et si} $Y$ \uline{a un faisceau ample}, \uline{en prenant pour} $V$ \uline{un} $\mathbf{P}^r_Y$).

\uline{Alors} $g\circ f$ \uline{est d'intersection complète}, \uline{lissifiable}, \uline{et on a un triangle distingué} :
\[
\begin{tikzcd}[column sep=large,row sep=large]
& \underline{L}^{X/Y}_\bullet \arrow[dl] \\
\mathbb Lf^*(\underline{L}^{Y/Z}_\bullet) \arrow[rr] && \underline{L}^{X/Z}_\bullet . \arrow[ul]
\end{tikzcd}
\]

Soit $V'$ un $Z$-schéma lisse factorisant $g$. On considère le diagramme commutatif :
\[
\begin{tikzcd}[column sep=large,row sep=large]
X \arrow[r,"i"] \arrow[dr,"f"'] & V \arrow[r,"j'"] \arrow[d,"p"] & V'' \arrow[d,"p''"]\\
& Y \arrow[r,"j"] \arrow[dr,"g"'] & V' \arrow[d,"p'"]\\
& & Z .
\end{tikzcd}
\]

Alors l'immersion $j'$ est régulière, car déduite de $j$ par un changement de base plat. Par suite $j'\circ i$ est une immersion régulière factorisant $g\circ f$, qui est donc lissifiable et d'intersection complète. De plus, on a la suite exacte des faisceaux conormaux :
\[
0\longrightarrow i^*(\underline{J}'/\underline{J}'^2)\longrightarrow \underline{K}/\underline{K}^2\longrightarrow \underline{I}/\underline{I}^2\longrightarrow0,
\]
en notant $\underline{I},\underline{J},\underline{J}',\underline{K}$ les faisceaux définissant respectivement les immersions $i,j,j',j'\circ i$. Or, $p''$ étant plat, $\underline{J}'=p''^*(\underline{J})$. La suite peut donc s'écrire :
\begin{equation}\tag{2.6.1}
0\longrightarrow f^*(\underline{J}/\underline{J}^2)\longrightarrow \underline{K}/\underline{K}^2\longrightarrow \underline{I}/\underline{I}^2\longrightarrow0.
\end{equation}

Par ailleurs, $V'$ et $V''$ étant lisses sur $Z$, on a la suite exacte :
\[
0\longrightarrow p''^*(\underline{\Omega}^1_{V'/Z})\longrightarrow\underline{\Omega}^1_{V''/Z}\longrightarrow\underline{\Omega}^1_{V''/V'}\longrightarrow0.
\]
